\makeabstract
    {
    Localizing Hermansky-Pudlak Syndrome complexes BLOC-1 and BLOC-2 during melanosome biogenesis
    }
    {
    Andrea J. Yan\textsuperscript{1}, Rachel M. Welles\textsuperscript{2,3}, Rachel Tocci\textsuperscript{3}, Chloe E. Snider\textsuperscript{3}, Michael S. Marks\textsuperscript{3,4,5}
    }
    {
    \textsuperscript{1} Amherst College, Amherst, MA \\
\textsuperscript{2} Cell and Molecular Biology Graduate Group, University of Pennsylvania, Philadelphia, PA\\
\textsuperscript{3} Department of Pathology and Laboratory Medicine, Children{\textquoteright}s Hospital of Philadelphia Research Institute, Philadelphia, PA\\
\textsuperscript{4} Department of Pathology and Laboratory Medicine, Perelman School of Medicine, Philadelphia, PA\\
\textsuperscript{5} Department of Physiology, Perelman School of Medicine, Philadelphia, PA
    }
    {
    Hermansky-Pudlak Syndrome (HPS) is a rare genetic disorder characterized principally by oculocutaneous albinism, with its severe visual impairment, and by excessive bleeding and bruising among other variable symptoms. HPS is caused by mutations in genes related to the biogenesis and function of lysosome-related organelles (LROs), which are cell type-specific organelles derived in part from the endolysosomal system. Among these LROs are melanosomes, which control the synthesis of the melanin that gives skin and eyes pigment. A key step of melanin synthesis is the maturation of non-pigmented melanosome precursors to fully pigmented melanosomes. Specifically, BLOC-1 (Biogenesis of Lysosome-related Organelles Complex 1) promotes the formation of tubules from endosomes that carry key cargo needed for melanin synthesis towards maturing melanosomes. BLOC-2 works with BLOC-1 during this transport process by promoting contact of the tubules with maturing melanosomes. While it is clear that BLOC-1 functions at the endosome membrane to initiate tubule formation, whether it also stabilizes the tubule and participates with BLOC-2 during tubule contact with melanosomes is not known.
    }
    {
    To visualize BLOC-1 in relation to tubules and BLOC-2, we created recombinant retroviruses that express a BLOC-1 subunit, muted, fused with a HaloTag protein. We are using this virus to infect unpigmented melanocytes that are missing the BLOC-1 muted subunit; thus, after expression of the missing BLOC-1 subunit, we expect the melanocytes to regain pigment.
    }
    {
    Cells lacking the muted subunit are normally not pigmented, and die off in the presence of hygromycin (hygro) selection. Both muted-rescue constructs (Muted and Muted-HaloTag) have surviving pigmented cells in the presence of hygro selection. All cell lines have living cells present, except for untransduced cells in the presence of hygro selection. Cells lacking the muted subunit are alive, but have no pigment. Thus, our construct works, restoring pigment as expected. Live cell fluorescence microscopy reveals little to no HaloDye signal in cells without the HaloTag fusion. There is a diffuse cytosolic HaloDye signal in muted-HaloTag rescues, which we believe to be due to cleavage of the HaloTag from the muted subunit.
    }
    {
    HaloTag does not interfere with BLOC-1 function as pigment is restored in melan-mu cells expressing the muted HaloTag fusion. We must run a Western Blot to investigate cleavage of HaloTag. From there, we can rescue with a double HaloTag construct, use protease inhibitors, and/or create a new muted KO cell line with a different parental cell line. After, we should visualize muted-HaloTag rescues again, hopefully seeing BLOC-1 puncta instead of a cytosolic HaloDye signal.
    }

    \newpage
    